%-------------------------------------------------------------------------------
%	TITRE DE SECTION
%-------------------------------------------------------------------------------
\cvsection{Expérience professionnelle}

%-------------------------------------------------------------------------------
%	CONTENU
%-------------------------------------------------------------------------------
\begin{cventries}

%---------------------------------------------------------
\cventry
  {Ingénieur Ruby on Rails senior}
  {Abbove (PaxFamilia) – Technologie financière patrimoniale}
  {Bruxelles, Belgique – Télétravail}
  {Juil. 2025 – Aujourd'hui (9 mois)}
  {
    \begin{cvitems}
      \item {Conçu et livré un système d'allocation géographique complet pour les portefeuilles d'investissement — modèle de données, services serveur, interfaces de programmation, graphiques multi-secteurs et formulaires — déployé progressivement par drapeaux de fonctionnalité pour des clients institutionnels (BNP Paribas).}
      \item {Construit un moteur de périodes de propriété temporelles pour les valorisations d'actifs, permettant le suivi de la propriété dans le temps avec des calculs de pourcentages adaptés au contexte sur une plateforme multi-locataire complexe de planification patrimoniale.}
      \item {Développé des services d'analyse de marchés privés et le suivi des valeurs de passifs, avec interpolation de séries temporelles pour les visualisations historiques.}
      \item {Maintenu et renforcé les intégrations bancaires tierces (Powens, Flanks) pour la synchronisation de portefeuilles, avec des suites de tests d'intégration complètes.}
      \item {Refactorisé systématiquement du code ancien : extraction de modules partagés, objets de service depuis des contrôleurs surchargés, et simplification du moteur de calcul de valorisation — plus de 32 000 lignes supprimées.}
      \item {Rédigé et exécuté des migrations de données multi-locataires avec des stratégies de déploiement par locataire, et corrigé des bugs critiques de valorisation (pourcentages de propriété, montants monétaires sur grands entiers).}
      \item {Étendu le moteur de rapports avec de nouveaux types de tableaux, colonnes et composants pour les vues de portefeuilles d'investissement et d'assurance.}
    \end{cvitems}
  }

%---------------------------------------------------------
\cventry
  {Ingénieur Ruby on Rails senior}
  {Climb (ex Tacotax) – Technologie financière}
  {Paris, France – Hybride}
  {Janv. 2023 – Aujourd'hui (2 ans 3 mois)}
  {
    \begin{cvitems}
      \item {Migré la base de code de Rails 6 vers Rails 7.2 et supprimé les composants obsolètes en lien avec un nouveau modèle économique.}
      \item {Construit des pipelines d'intégration et de déploiement continus avec des scripts de déploiement sur mesure pour automatiser et fluidifier les mises en production.}
      \item {Implémenté des intégrations d'interfaces de programmation entre entreprises critiques pour le métier et intégré Stripe pour le traitement sécurisé des paiements.}
      \item {Connecté des interfaces de programmation partenaires externes à Salesforce pour synchroniser les données de contrats clients en temps réel.}
      \item {Relié les données Salesforce à l'interface pour une expérience utilisateur fluide et riche en données.}
      \item {Réduit les erreurs côté interface et côté serveur à zéro dans Rollbar grâce à un débogage et une surveillance proactifs.}
      \item {Amélioré le temps de réponse des interfaces de programmation et les performances système via des optimisations ciblées de code et de requêtes.}
      \item {Contribué à ClimbGPT, un assistant d'intelligence artificielle pour le conseil fiscal et patrimonial utilisant l'interface de programmation OpenAI.}
    \end{cvitems}
  }

%---------------------------------------------------------
\cventry
  {Ingénieur polyvalent Ruby on Rails et React}
  {Doctolib – Technologie médicale}
  {Paris, France – Hybride}
  {Juin 2020 – Déc. 2022 (2 ans 6 mois)}
  {
    \begin{cvitems}
      \item {Construit l'application de prescription utilisée par tous les médecins en France, au sein d'une équipe produit pluridisciplinaire.}
      \item {Appliqué les bonnes pratiques : programmation en binôme et en mob, développement piloté par les tests et revues de code.}
      \item {Maintenu une couverture de tests élevée avec Minitest et Capybara pour garantir la fiabilité des fonctionnalités.}
      \item {Collaboré avec les chefs de produit pour définir les spécifications, jalons et documentation technique.}
      \item {Formé de nouveaux ingénieurs sur Rails et React ; accompagné l'intégration de nouveaux membres dans l'équipe.}
      \item {Participé à l'astreinte développement avec des déploiements en production quotidiens (3 fois par jour).}
    \end{cvitems}
  }

%---------------------------------------------------------
\cventry
  {Développeur React/Next.js et fondateur}
  {Lingu.Africa – Technologie de contenu}
  {Télétravail}
  {Mars 2021 – Aujourd'hui (3 ans et plus)}
  {
    \begin{cvitems}
      \item {Construit une plateforme de commerce en ligne multilingue pour plus de 180 livres pour enfants auto-édités.}
      \item {Créé un système de blog dynamique avec conversion Markdown vers HTML et optimisation du référencement.}
      \item {Automatisé les mises à jour de catalogue, l'analyse de données et la publication de blogs générés par intelligence artificielle (articles de plus de 21 000 mots avec images dynamiques).}
      \item {Obtenu un bon référencement grâce à des pages programmatiques, métadonnées, plans de site et balises pour les réseaux sociaux.}
      \item {Géré l'architecture, la conception et le déploiement avec Next.js et TypeScript.}
    \end{cvitems}
  }

%---------------------------------------------------------
\cventry
  {Ingénieur logiciel Ruby on Rails}
  {Workelo – Technologie RH}
  {Paris, France – Sur site}
  {Juil. 2019 – Juin 2020 (11 mois)}
  {
    \begin{cvitems}
      \item {Refactorisé et fait évoluer le système de messagerie interne pour améliorer la fiabilité.}
      \item {Construit des pipelines d'intégration et de déploiement continus avec GitHub Actions et CircleCI.}
      \item {Introduit les tests unitaires et d'intégration, et formé l'équipe aux pratiques de test.}
      \item {Maintenu la qualité du code avec RuboCop, SimpleCov, CodeClimate et Brakeman.}
    \end{cvitems}
  }

%---------------------------------------------------------
\cventry
  {Formateur Ruby on Rails}
  {Le Wagon – Technologie éducative}
  {Paris, France – Sur site}
  {Janv. 2019 – Juin 2019 (5 mois)}
  {
    \begin{cvitems}
      \item {Enseigné Ruby, programmation orientée objet et les fondamentaux Rails à des étudiants en développement web à temps plein.}
      \item {Animé des sessions pratiques et accompagné des équipes projet sur les bonnes pratiques.}
    \end{cvitems}
  }

%---------------------------------------------------------
\cventry
  {Ingénieur Ruby on Rails}
  {Joinly – Technologie sportive}
  {Paris, France – Sur site}
  {Janv. 2018 – Janv. 2019 (1 an)}
  {
    \begin{cvitems}
      \item {Construit une plateforme logicielle en tant que service de gestion des paiements d'arbitres avec intégration de l'interface de programmation S-Money.}
      \item {Implémenté des pipelines d'intégration continue avec CircleCI et GitHub Actions.}
      \item {Maintenu la qualité du code avec des outils d'analyse de sécurité et de performance.}
    \end{cvitems}
  }

%---------------------------------------------------------
\cventry
  {Ingénieur logiciel et réseau}
  {Covage – Télécommunications}
  {Paris, France – Sur site}
  {Fév. 2015 – Sept. 2017 (2 ans 7 mois)}
  {
    \begin{cvitems}
      \item {Automatisé les déploiements de services réseau et les audits de configuration.}
      \item {Implémenté l'infrastructure en tant que code et réduit les opérations manuelles.}
      \item {Géré le carnet de produit, encadré des stagiaires et piloté le développement de fonctionnalités.}
    \end{cvitems}
  }

%---------------------------------------------------------
\end{cventries}
